% Cas d'utilisation « Gérer les activités planifiées » — inclus sous la sous-section 4.4.4
\begin{table}[htbp]
  \centering
  \small
  \caption{Description textuelle du cas d'utilisation : «~Gérer les activités planifiées~»}
  \label{tab:cu-gerer-activites-planifiees}
  \PFETableSetup
  \PFETableRowColors
  \PFETableArrayStretch
  \PFETableTabularXBegin
  \begin{tabularx}{\PFETableBlockWidth}{@{}>{\raggedright\arraybackslash}p{3.2cm} >{\raggedright\arraybackslash}X@{}}
    \tableHeaderStyle \tableHeaderCell{Champ} & \tableHeaderCell{Contenu} \\
    \textbf{Titre} & Gérer les activités planifiées \\
    \textbf{Acteurs} & Utilisateur site (saisie et mise à jour), utilisateur corporate (contrôle selon les droits du workflow). \\
    \textbf{Objectif} &
      Décliner le plan annuel en lignes d'activités prévisionnelles (thématique, budget, indicateurs) et maintenir ces lignes tout au long de la phase de préparation jusqu'au verrouillage lié au statut du plan. \\
    \textbf{Préconditions} &
      Un plan CSR existe pour le site. L'utilisateur site a accès au plan. Le statut du plan et de l'activité autorise les modifications lorsque les règles métier le permettent. \\
    \textbf{Postconditions} &
      Les activités planifiées sont persistées et visibles dans la vue détail du plan ; elles servent de référence pour le suivi réalisé ultérieur. \\
    \textbf{Scénario nominal} &
      L'utilisateur site ouvre le détail d'un plan, ajoute ou modifie des lignes d'activités, enregistre. Le système valide les champs obligatoires et met à jour les totaux ou compteurs affichés. \\
    \textbf{Scénario alternatif} &
      \textbf{Plan ou activité non éditable :} après soumission ou validation, certaines actions sont masquées ; le système informe l'utilisateur des règles de modification ou renvoie vers une demande de changement si applicable.\par\medskip
      \textbf{Hors plan :} création d'une activité hors plan avec un circuit d'approbation distinct lorsque l'action n'entre pas dans le périmètre strict du plan validé. \\
  \end{tabularx}%
  \PFETableTabularXEnd
\end{table}
